% answers_short.tex —— 2013 年·陕西卷·理科 · 简短答案速查
% 由 tools/make_short_answers.py 自 answers.tex 的 \gksolhead/\gkanswer 提取，
% 并与 2013/2013_数学真题_陕西理_答案要点.md 逐题比对后生成，勿手改（改源 MD 后重跑）。

\begin{gknotice}
\textbf{使用说明}：本节只给最终答案便于快速核对，完整推导见解析版。答案为编写者独立推导，\textbf{非官方标准答案}。
\end{gknotice}

\gkpart{选择题}
\begin{center}\begin{tabular}{ccccccccccc}
\toprule
题号 & 1 & 2 & 3 & 4 & 5 & 6 & 7 & 8 & 9 & 10 \\
\midrule
答案 & D & C & C & B & A & D & B & A & C & D \\
\bottomrule
\end{tabular}\end{center}

\gkpart{填空题}
\begin{itemize}[leftmargin=2.6em]
  \item[\textbf{11.}] $9$。$a=4$，$e=\dfrac ca=\dfrac54$ 得 $c=5$，故 $b^2=c^2-a^2=9$，即 $m=9$。
  \item[\textbf{12.}] $\dfrac{\pi}{3}$。三视图还原为底面半径 $1$、高 $2$ 的半圆锥，体积 $V=\dfrac13\cdot\dfrac{\pi}{2}\cdot2=\dfrac\pi3$。
  \item[\textbf{13.}] $-4$。区域为 $|x-1|\leqslant y\leqslant2$，故 $x\geqslant-1$，$y\leqslant2$，从而 $2x-y\geqslant2(-1)-2=-4$，在 $(-1,2)$ 处取到。
  \item[\textbf{14.}] $1^2-2^2+3^2-\cdots+(-1)^{n+1}n^2=(-1)^{n+1}\dfrac{n(n+1)}2$。左边为前 $n$ 个平方的交错和，按 $n$ 的奇偶归纳即得。
  \item[\textbf{15.}] $2$。展开得 $(am+bn)(bm+an)=ab(m^2+n^2)+2(a^2+b^2)=ab(m^2+n^2-4)+2\geqslant2$，当 $m=n=\sqrt2$ 时取等号。
  \item[\textbf{16.}] $\sqrt6$。由 $PE\parallel BC$ 及 $A$，$B$，$C$，$D$ 共圆得 $\triangle PAE\backsim\triangle PED$，故 $PE^2=PA\cdot PD=3\times2=6$。
  \item[\textbf{17.}] $\begin{cases}x=\cos^2\theta,\\ y=\sin\theta\cos\theta\end{cases}$（$\theta$ 为参数，$0\leqslant\theta<\pi$）。将 $y=x\tan\theta$ 代入 $x^2+y^2-x=0$，舍去原点得 $x=\cos^2\theta$。
\end{itemize}

\gkpart{解答题}
\begin{itemize}[leftmargin=2.6em]
  \item[\textbf{18.}] (1) $T=\pi$。$f(x)=\sin\left(2x-\dfrac\pi6\right)$，故最小正周期为 $\pi$；(2) 最大值为 $1$，最小值为 $-\dfrac12$。在 $\left[0,\dfrac\pi2\right]$ 上 $2x-\dfrac\pi6\in\left[-\dfrac\pi6,\dfrac{5\pi}6\right]$
  \item[\textbf{19.}] (1) $S_n=\begin{cases}na_1,&q=1,\\[4pt]\dfrac{a_1(1-q^n)}{1-q},&q\neq1.\end{cases}$ 由错位相减可得；(2) 数列 $\{a_n+1\}$ 不是等比数列，见证明。若它为等比数列，由 $b_2^2=b_1b_3$ 可推出 $q=1$，与 $q\neq1$ 矛盾
  \item[\textbf{20.}] (1) $A_1C\perp$ 平面 $BB_1D_1D$，见证明；(2) $\theta=\dfrac\pi3$。两平面的法向量可分别取 $(1,0,1)$ 与 $(0,1,-1)$，夹角的余弦为 $\dfrac12$
  \item[\textbf{21.}] (1) $\dfrac{4}{15}$。$P(\text{\text{甲选} }3)=\dfrac23$，$P(\text{\text{乙不选} }3)=\dfrac{\mathrm C_4^3}{\mathrm C_5^3}=\dfrac25$，故所求为 $\dfrac23\times\dfrac25=\dfrac4{15}$；(2) $P(X=0)=\dfrac4{75}$，$P(X=1)=\dfrac4{15}$，$P(X=2)=\dfrac{11}{25}$，$P(X=3)=\dfrac6{25}$；$E(X)=\dfrac{28}{15}$。由 $X=X_{\text{\text{甲}}}+X_{\text{\text{乙}}}+X_{\text{\text{丙}}}$ 及独立性计算
  \item[\textbf{22.}] (1) $y^2=8x$。设圆心 $(x,y)$，半径 $r$，由 $r^2=(x-4)^2+y^2$ 与 $r^2=x^2+16$ 消去 $r$ 即得；(2) 直线 $l$ 过定点 $(1,0)$，见证明。设 $P(2u^2,4u)$，$Q(2v^2,4v)$，由角平分线条件得 $uv=-\dfrac12$，弦 $PQ$ 的方程恒过 $(1,0)$
  \item[\textbf{23.}] (1) $k=\dfrac1{\mathrm e^2}$。反函数为 $y=\ln x$，设切点 $x_0$，由斜率与截距条件得 $\ln x_0=2$；(2) 当 $0<m<\dfrac{\mathrm e^2}{4}$ 时无公共点；当 $m=\dfrac{\mathrm e^2}{4}$ 时恰有 $1$ 个公共点；当 $m>\dfrac{\mathrm e^2}{4}$ 时恰有 $2$ 个公共点。考查 $\dfrac{\mathrm e^x}{x^2}$ 在 $x>0$ 上的最小值；(3) $\dfrac{f(a)+f(b)}2>\dfrac{f(b)-f(a)}{b-a}$，见证明。令 $t=\dfrac{b-a}2>0$，用导数可证 $t(\mathrm e^t+\mathrm e^{-t})>\mathrm e^t-\mathrm e^{-t}$
\end{itemize}
